\documentclass{article}

\usepackage{amsfonts}

\begin{document}

\textbf{1.2 The Logistic Population Model}

\hfill

The differential equation $x'= ax$ above can be considered a simplistic model
of population growth when a > 0. The quantity x(t ) measures the population
of some species at time t . The assumption that leads to the differential equation
is that the rate of growth of the population (namely, dx/dt ) is directly
proportional to the size of the population. Of course, this naive assumption
omits many circumstances that govern actual population growth, including,
for example, the fact that actual populations cannot increase with bound.

To take this restriction into account, we can make the following further
assumptions about the population model:

\begin{enumerate}
\item If the population is small, the growth rate is nearly directly proportional
to the size of the population;
\item but if the population grows too large, the growth rate becomes negative.
\end{enumerate}

One differential equation that satisfies these assumptions is the logistic
population growth model. This differential equation is

\[
x' = ax (1 - \frac{x}{N}
\]

Here a and N are positive parameters: a gives the rate of population growth
when x is small, while N represents a sort of “ideal” population or “carrying
capacity.” Note that if x is small, the differential equation is essentially $x' = ax$
[since the term $1-(x/N) \approx 1$], but if $x > N$, then $x'< 0$. Thus this simple
equation satisfies the above assumptions. We should add here that there are
many other differential equations that correspond to these assumptions; our
choice is perhaps the simplest.

Without loss of generality we will assume that N = 1. That is, we will
choose units so that the carrying capacity is exactly 1 unit of population, and
x(t) therefore represents the fraction of the ideal population present at time t .

Therefore the logistic equation reduces to

\[
x' = f_a(x) = ax(1 - x).
\]

This is an example of a first-order, autonomous, nonlinear differential equation.
It is first order since only the first derivative of x appears in the equation.
It is autonomous since the right-hand side of the equation depends on x alone,
not on time t . And it is nonlinear since $f_a(x)$ is a nonlinear function of x. The
previous example, $x'= ax$, is a first-order, autonomous, linear differential
equation.

The solution of the logistic differential equation is easily found by the triedand-true calculus method of separation and integration:

\[
\int \frac{dx}{x(1-x)} = \int a dt
\]

The method of partial fractions allows us to rewrite the left integral as

\[
\int (\frac{1}{x} + \frac{1}{1-x}) dx
\]

Integrating both sides and then solving for x yields

\[
x(t) = \frac{K e^{at}}{1+ Ke^{at}}
\]

where K is the arbitrary constant that arises from integration. Evaluating this
expression at t = 0 and solving for K gives

\[
K = \frac{x(0)}{1 - x(0)}
\]

Using this, we may rewrite this solution as

\[
\frac{x(0)e^{at}}{1 - x(0) + x(0)e^{at}} .
\]


So this solution is valid for any initial population x(0). When x(0) = 1,
we have an equilibrium solution, since x(t) reduces to $x(t) \equiv 1$. Similarly,
$x(t) \equiv 0$ is also an equilibrium solution.

Thus we have “existence” of solutions for the logistic differential equation.
We have no guarantee that these are all of the solutions to this equation at this
stage; we will return to this issue when we discuss the existence and uniqueness
problem for differential equations in Chapter 7.

To get a qualitative feeling for the behavior of solutions, we sketch the
slope field for this equation. The right-hand side of the differential equation
determines the slope of the graph of any solution at each time t . Hence we
may plot little slope lines in the tx–plane as in Figure 1.3, with the slope of the
line at (t,x) given by the quantity $ax(1 - x)$. Our solutions must therefore
have graphs that are everywhere tangent to this slope field. From these graphs,
we see immediately that, in agreement with our assumptions, all solutions
for which $x(0) > 0$ tend to the ideal population $x(t) \equiv 1$. For $x(0) < 0$,
solutions tend to $- \infty$, although these solutions are irrelevant in the context
of a population model.

Note that we can also read this behavior from the graph of the function
$f_a(x) = ax(1 - x)$. This graph, displayed in Figure 1.4, crosses the x-axis at
the two points x = 0 and x = 1, so these represent our equilibrium points.
When 0 < x < 1, we have f (x) > 0. Hence slopes are positive at any (t , x)
with 0 < x < 1, and so solutions must increase in this region. When x < 0 or
x > 1, we have f (x) < 0 and so solutions must decrease, as we see in both the
solution graphs and the phase lines in Figure 1.3.
We may read off the fact that x = 0 is a source and x = 1 is a sink from
the graph of f in similar fashion. Near 0, we have f (x) > 0 if x > 0, so slopes
are positive and solutions increase, but if x < 0, then f (x) < 0, so slopes are
negative and solutions decrease. Thus nearby solutions move away from 0 and
so 0 is a source. Similarly, 1 is a sink.

We may also determine this information analytically. We have $f'_a(x)=a-2ax$ so that $f'_a(0) = a > 0$ and $f'_a(1) = -a < 0$. Since $f'_a(0) > 0$, slopes must increase through the value 0 as x passes through 0. That is, slopes are
negative below x = 0 and positive above x = 0. Hence solutions must tend
away from x = 0. In similar fashion, $f'_a(1) < 0$ forces solutions to tend toward
x = 1, making this equilibrium point a sink. We will encounter many such
“derivative tests” like this that predict the qualitative behavior near equilibria
in subsequent chapters.

\textbf{Example.} As a further illustration of these qualitative ideas, consider the
differential equation

\[
x' = g (x) = x - x^3.
\]


There are three equilibrium points, at $x = 0, \pm 1$. Since $g'(x) = 1 - 3x^2$, we have $g'(0) = 1$, so the equilibrium point 0 is a source. Also, $g'(\pm 1) = -2$, so the equilibrium points at $\pm 1$ are both sinks. Between these equilibria, the sign of the slope field of this equation is nonzero. From this information we can immediately display the phase line, which is shown in Figure 1.5.  

\hfill

\textbf{1.3 Constant Harvesting and
Bifurcations}

\hfill

Now let’s modify the logistic model to take into account harvesting of the population.
Suppose that the population obeys the logistic assumptions with the

parameter a = 1, but is also harvested at the constant rate h. The differential
equation becomes

\[
x  = x(1 - x) - h
\]

where $h \geq 0$ is a new parameter.

Rather than solving this equation explicitly (which can be done — see
Exercise 6 at the end of this chapter), we use the graphs of the functions

\[
f_h(x) = x(1 - x) - h
\]

to “read off ” the qualitative behavior of solutions. In Figure 1.6 we display
the graph of $f_h$ in three different cases: 0 < h < 1/4, h = 1/4, and h > 1/4.
It is straightforward to check that fh has two roots when $0 \leq h < 1/4$, one
root when h = 1/4, and no roots if $h > 1/4$, as illustrated in the graphs. As a
consequence, the differential equation has two equilibrium points $x_l$  and $x_r$
with $0 \leq x  < x_r$ when $0 < h < 1/4$. It is also easy to check that $f'_h (x) > 0$, so
that $x_l$  is a source, and $f'_h (x_r) < 0$ so that $x_r$ is a sink.


As h passes through h = 1/4, we encounter another example of a bifurcation.
The two equilibria $x_l$  and $x_r$ coalesce as h increases through 1/4 and then
disappear when h > 1/4. Moreover, when h > 1/4, we have $f_h(x) < 0$ for all
x. Mathematically, this means that all solutions of the differential equation
decrease to $-\infty$ as time goes on.


We record this visually in the bifurcation diagram. In this diagram we plot
the parameter h horizontally. Over each h-value we plot the corresponding
phase line. The curve in this picture represents the equilibrium points for each
value of h. This gives another view of the sink and source merging into a
single equilibrium point and then disappearing as h passes through 1/4 (see
Figure 1.7).


Ecologically, this bifurcation corresponds to a disaster for the species under
study. For rates of harvesting 1/4 or lower, the population persists, provided

the initial population is sufficiently large ($x(0) \geq x$ ). But a very small change
in the rate of harvesting when h = 1/4 leads to a major change in the fate of
the population: At any rate of harvesting h > 1/4, the species becomes extinct.
This phenomenon highlights the importance of detecting bifurcations in
families of differential equations, a procedure that we will encounter many
times in later chapters. We should also mention that, despite the simplicity of
this population model, the prediction that small changes in harvesting rates
can lead to disastrous changes in population has been observed many times in
real situations on earth.
Example. As another example of a bifurcation, consider the family of
differential equations
x  = ga(x) = x2 − ax = x(x − a)
which depends on a parameter a. The equilibrium points are given by x = 0
and x = a. We compute g  
a(0) = -a, so 0 is a sink if a > 0 and a source if
a < 0. Similarly, g  
a(a) = a, so x = a is a sink if a < 0 and a source if a > 0.
We have a bifurcation at a = 0 since there is only one equilibrium point when
a = 0. Moreover, the equilibrium point at 0 changes from a source to a sink
as a increases through 0. Similarly, the equilibrium at x = a changes from a
sink to a source as a passes through 0. The bifurcation diagram for this family
is depicted in Figure 1.8.  

\hfill

\textbf{1.4 Periodic Harvesting and
Periodic Solutions}

\hfill

Now let’s change our assumptions about the logistic model to reflect the
fact that harvesting does not always occur at a constant rate. For example,

populations of many species of fish are harvested at a higher rate in warmer
seasons than in colder months. So we assume that the population is harvested
at a periodic rate. One such model is then

\[
x'  = f (t , x) = ax(1 - x) - h(1 + \sin(2 \pi t ))
\]


where again a and h are positive parameters. Thus the harvesting reaches a
maximum rate -2h at time t = 14
+ n where n is an integer (representing the
year), and the harvesting reaches its minimum value 0 when t = 34
+n, exactly
one-half year later. Note that this differential equation now depends explicitly
on time; this is an example of a nonautonomous differential equation. As in
the autonomous case, a solution x(t ) of this equation must satisfy x (t ) =
f (t , x(t )) for all t . Also, this differential equation is no longer separable, so we
cannot generate an analytic formula for its solution using the usual methods
from calculus. Thus we are forced to take a more qualitative approach.
To describe the fate of the population in this case, we first note that the righthand
side of the differential equation is periodic with period 1 in the time
variable. That is, f (t + 1, x) = f (t , x). This fact simplifies somewhat the
problem of finding solutions. Suppose that we know the solution of all initial
value problems, not for all times, but only for 0 ≤ t ≤ 1. Then in fact we
know the solutions for all time. For example, suppose x1(t ) is the solution that
is defined for 0 ≤ t ≤ 1 and satisfies x1(0) = x0. Suppose that x2(t) is the
solution that satisfies x2(0) = x1(1). Then we may extend the solution x1 by
defining x1(t + 1) = x2(t ) for $0 \leq t \leq 1$. The extended function is a solution
since we have
x 
1(t + 1) = x 
2(t ) = f (t , x2(t ))
= f (t + 1, x1(t + 1)).

Thus if we know the behavior of all solutions in the interval 0 ≤ t ≤ 1, then
we can extrapolate in similar fashion to all time intervals and thereby know
the behavior of solutions for all time.
Secondly, suppose that we know the value at time t = 1 of the solution
satisfying any initial condition x(0) = x0. Then, to each such initial condition
x0, we can associate the value x(1) of the solution x(t ) that satisfies x(0) = x0.
This gives us a function p(x0) = x(1). If we compose this function with itself,
we derive the value of the solution through x0 at time 2; that is, p(p(x0)) =
x(2). If we compose this function with itself n times, then we can compute
the value of the solution curve at time n and hence we know the fate of the
solution curve.
The function p is called a Poincaré map for this differential equation. Having
such a function allows us to move from the realm of continuous dynamical
systems (differential equations) to the often easier-to-understand realm of
discrete dynamical systems (iterated functions). For example, suppose that
we know that p(x0) = x0 for some initial condition x0. That is, x0 is a fixed
point for the function p. Then from our previous observations, we know that
x(n) = x0 for each integer n. Moreover, for each time t with 0 < t < 1, we
also have x(t ) = x(t + 1) and hence x(t + n) = x(t ) for each integer n.
That is, the solution satisfying the initial condition x(0) = x0 is a periodic
function of t with period 1. Such solutions are called periodic solutions of the
differential equation. In Figure 1.10, we have displayed several solutions of the
logistic equation with periodic harvesting. Note that the solution satisfying
the initial condition x(0) = x0 is a periodic solution, and we have x0 =
p(x0) = p(p(x0)) . . .. Similarly, the solution satisfying the initial condition
x(0) = ˆx0 also appears to be a periodic solution, so we should have p(ˆx0) = ˆx0.
Unfortunately, it is usually the case that computing a Poincaré map for a
differential equation is impossible, but for the logistic equation with periodic
harvesting we get lucky.

\hfill

\textbf{1.5 Computing the Poincaré Map}

\hfill

Before computing the Poincaré map for this equation, we introduce some
important terminology. To emphasize the dependence of a solution on the
initial value $x_0$, we will denote the corresponding solution by $\phi(t,x_0)$. This
function $\phi: \mathbb{R} \times \mathbb{R} \to  \mathbb{R}$ is called the flow associated to the differential
equation. If we hold the variable x0 fixed, then the function

\[
t \to \phi(t, x_0)
\]

is just an alternative expression for the solution of the differential equation
satisfying the initial condition x0. Sometimes we write this function as $\phi_t(x_0)$.

\hfill

\textbf{Example.} For our first example, x' = ax, the flow is given by

\[
\phi(t , x_0) = x_0e^{at}.
\]

For the logistic equation (without harvesting), the flow is

\[
\phi(t, x_0) = \frac{x(0)e^{at}}{1 - x(0) + x(0)e^{at}
\]

Now we return to the logistic differential equation with periodic harvesting

\[
x'  = f (t , x) = ax(1 - x) - h(1 + \sin(2 \pi t )).
\]


The solution satisfying the initial condition $x(0) = x_0$ is given by $t \to \phi(t, x_0)$.
While we do not have a formula for this expression, we do know that, by the
fundamental theorem of calculus, this solution satisfies

\[
\phi(t, x_0) = x_0 + \int_0^t f(s, \phi(s, x_0)) ds
\]

and $\phi(0, x_0) = x_0$.

If we differentiate this solution with respect to x0, we obtain, using the chain
rule:

\[
\frac{\partial \phi}{\partial x_0}(t, x_0) = 1 + \int_0^t \frac{\partial f}{\partial x_0} (s, \phi(s, x_0)) \dot \frac{\partial \phi}{\partial x_0} (s, x_0) ds
\]

Now let

\[
z(t) =  \frac{\partial \phi}{\partial x_0}(t, x_0)
\]

Differentiating z with respect to t, we find

\[
z'(t ) = \frac{\partial f}{\partial x_0}(t, \phi(t, x_0)) \dot \frac{\partial \phi}{\partial x_0}(t, x_0)
\]
\[
= \frac{\partial f}{\partial x_0} (t, \phi/t, x_0)) \dot z(t)
\]

Again, we do not know φ(t , x0) explicitly, but this equation does tell us that
z(t ) solves the differential equation
z (t ) = ∂f
∂x0
(t , φ(t , x0)) z(t )
with z(0) = 1. Consequently, via separation of variables, we may compute
that the solution of this equation is
z(t ) = exp
  t
0
∂f
∂x0
(s, φ(s, x0)) ds

and so we find
∂φ
∂x0
(1, x0) = exp
  1
0
∂f
∂x0
(s, φ(s, x0)) ds.
Since p(x0) = φ(1, x0), we have determined the derivative p (x0) of the
Poincaré map. Note that p (x0) > 0. Therefore p is an increasing function.
Differentiating once more, we find
p  (x0) = p (x0)
   1
0
∂2f
∂x0∂x0
(s, φ(s, x0)) · exp
   s
0
∂f
∂x0
(u, φ(u, x0)) du
 
ds
 
,
which looks pretty intimidating. However, since
f (t , x0) = ax0(1 − x0) − h(1 + sin(2πt )),
we have
∂2f
∂x0∂x0
≡ −2a.
Thus we know in addition that p  (x0) < 0. Consequently, the graph of the
Poincaré map is concave down. This implies that the graph of p can cross the
diagonal line y = x at most two times. That is, there can be at most two values
of x for which p(x) = x. Therefore the Poincaré map has at most two fixed
points. These fixed points yield periodic solutions of the original differential
equation. These are solutions that satisfy x(t +1) = x(t ) for all t . Another way
to say this is that the flow φ(t , x0) is a periodic function in t with period 1 when
the initial condition x0 is one of the fixed points. We saw these two solutions
in the particular case when h = 0. 8 in Figure 1.10. In Figure 1.11, we again
see two solutions that appear to be periodic. Note that one of these solutions
appears to attract all nearby solutions, while the other appears to repel them.
We will return to these concepts often and make them more precise later in
the book.
Recall that the differential equation also depends on the harvesting parameter
h. For small values of h there will be two fixed points such as shown in
Figure 1.11. Differentiating f with respect to h, we find
∂f
∂h
(t , x0) = −(1 + sin 2πt )
Hence ∂f /∂h < 0 (except when t = 3/4). This implies that the slopes of the
slope field lines at each point (t , x0) decrease as h increases. As a consequence,
the values of the Poincaré map also decrease as h increases. Hence there is a
unique value h∗ for which the Poincaré map has exactly one fixed point. For
h > h∗, there are no fixed points for p and so p(x0) < x0 for all initial values.
It then follows that the population again dies out.

\hfill

\textbf{1.6 Exploration: A Two-Parameter
Family}

\hfill

Consider the family of differential equations

\[
x' = f_{a,b}(x) = ax − x^3 − b
\]

which depends on two parameters, a and b. The goal of this exploration is
to combine all of the ideas in this chapter to put together a complete picture
of the two-dimensional parameter plane (the ab–plane) for this differential
equation. Feel free to use a computer to experiment with this differential
equation at first, but then try to verify your observations rigorously.

\begin{enumerate}
\item First fix a = 1. Use the graph of f1,b to construct the bifurcation diagram
for this family of differential equations depending on b.
\item Repeat the previous step for a = 0 and then for a = −1.
\item What does the bifurcation diagram look like for other values of a?
\item Now fix b and use the graph to construct the bifurcation diagram for this
family, which this time depends on a.
\item In the ab–plane, sketch the regions where the corresponding differential
equation has different numbers of equilibrium points, including a sketch
of the boundary between these regions.
\item Describe, using phase lines and the graph of fa,b(x), the bifurcations that
occur as the parameters pass through this boundary.
\item Describe in detail the bifurcations that occur at a = b = 0 as a and/or b
vary.

\item Consider the differential equation x  = x − x3 − b sin(2πt ) where |b| is
small. What can you say about solutions of this equation? Are there any
periodic solutions?
\item Experimentally, what happens as |b| increases? Do you observe any
bifurcations? Explain what you observe.
\end{enumerate}

\hfill

\textbf{EXERCISES}

\hfill

\begin{enumerate}
\item Find the general solution of the differential equation x  = ax + 3 where
a is a parameter. What are the equilibrium points for this equation? For
which values of a are the equilibria sinks? For which are they sources?

\item For each of the following differential equations, find all equilibrium
solutions and determine if they are sinks, sources, or neither. Also, sketch
the phase line.

\begin{enumerate}
\item  x  = x3 − 3x
\item  x  = x4 − x2
\item  x  = cos x
\item  x  = sin2 x
\item  x  = |1 − x2|
\end{enumerate}

\item Each of the following families of differential equations depends on a
parameter a. Sketch the corresponding bifurcation diagrams.

\begin{enumerate}
\item  x  = x2 − ax
\item x  = x3 − ax
\item  x  = x3 − x + a
\end{enumerate}

\item Consider the function f (x) whose graph is displayed in Figure 1.12.

\begin{enumerate}
\item Sketch the phase line corresponding to the differential equation x  =
f (x).
\item Let ga(x) = f (x)+a. Sketch the bifurcation diagram corresponding
to the family of differential equations x  = ga(x).
\item Describe the different bifurcations that occur in this family.
\end{enumerate}

\item Consider the family of differential equations

\[
x'  = ax + \sin x
\]

where a is a parameter.

\begin{enumerate}
\item Sketch the phase line when a = 0.
\item Use the graphs of ax and sin x to determine the qualitative behavior
of all of the bifurcations that occur as a increases from −1 to 1.
\item Sketch the bifurcation diagram for this family of differential
equations.
\end{enumerate}

\item Find the general solution of the logistic differential equation with
constant harvesting

\[
x'  = x(1 - x) - h
\]

for all values of the parameter h > 0.

\item Consider the nonautonomous differential equation

\[
 x'=\left\{\begin{matrix}
x-4 \text{ if } t<5 \\
2-x \text{ if } t \geq 5
\end{matrix}\right.
\]

\begin{enumerate}
\item Find a solution of this equation satisfying x(0) = 4. Describe the
qualitative behavior of this solution.
\item Find a solution of this equation satisfying x(0) = 3. Describe the
qualitative behavior of this solution.
\item Describe the qualitative behavior of any solution of this system as
$t \to \infty$.
\end{enumerate}

\item Consider a first-order linear equation of the form x  = ax + f (t ) where
$a \in R$. Let y(t ) be any solution of this equation. Prove that the general
solution is y(t ) + c exp(at ) where c ∈ R is arbitrary.

\item Consider a first-order, linear, nonautonomous equation of the form
x (t ) = a(t )x.

\begin{enumerate}
\item Find a formula involving integrals for the solution of this system.
\item Prove that your formula gives the general solution of this system.
\end{enumerate}

\item Consider the differential equation x  = x + cos t .

\begin{enumerate}
\item Find the general solution of this equation.
\item Prove that there is a unique periodic solution for this equation.
\item Compute the Poincaré map p: {t = 0} → {t = 2π} for this equation
and use this to verify again that there is a unique periodic
solution.
\end{enumerate}

\item First-order differential equations need not have solutions that are defined
for all times.

\begin{enumerate}
\item Find the general solution of the equation x  = x2.
\item Discuss the domains over which each solution is defined.
\item Give an example of a differential equation for which the solution
satisfying x(0) = 0 is defined only for −1 < t < 1.
\end{enumerate}

\item First-order differential equations need not have unique solutions satisfying
a given initial condition.

\begin{enumerate}
\item Prove that there are infinitely many different solutions of the
differential equations x  = x1/3 satisfying x(0) = 0.
\item Discuss the corresponding situation that occurs for x  = x/t ,
x(0) = x0.
\item Discuss the situation that occurs for x  = x/t 2, x(0) = 0.
\end{enumerate}

\item Let x  = f (x) be an autonomous first-order differential equation with
an equilibrium point at x0.

\begin{enumerate}
\item Suppose f  (x0) = 0. What can you say about the behavior of
solutions near x0? Give examples.
\item Suppose f  (x0) = 0 and f   (x0)  = 0. What can you now say?
\item Suppose f  (x0) = f   (x0) = 0 but f    (x0)  = 0. What can you now
say?
\end{enumerate}

\item Consider the first-order nonautonomous equation x' = p(t)x where
p(t ) is differentiable and periodic with period T. Prove that all solutions
of this equation are periodic with period T if and only if

\[
\int_0^T p(s) ds = 0
\]

\item Consider the differential equation x  = f (t , x) where f (t , x) is continuously
differentiable in t and x. Suppose that

\[
f (t + T, x) = f (t , x)
\]

for all t . Suppose there are constants p, q such that

\[
f (t , p) > 0, f (t , q) < 0
\]

for all t . Prove that there is a periodic solution x(t ) for this equation with
p < x(0) < q.
16. Consider the differential equation x  = x2 − 1 − cos(t ). What can be
said about the existence of periodic solutions for this equation?



\end{document}



